% FILE: 03_architecture_and_primitives.tex
% Project: standards
% Thesis Role: public-standards-compliance-and-gravity-field-mapping

\section{Architecture and Primitives}
\label{sec:standards-arch}

\subsection{Core Objects}
\label{sec:standards-arch-objects}

The standards project is organized around four classes of core objects:

\paragraph{Standard placement files.}
Each of the 16 standards receives a dedicated placement file that defines: (1) the
ontological mapping from the standard's data elements to ledger assertions; (2) the
type-law surface expressed as constraints on the \texttt{Evidence<T, State, Witness>}
lattice; (3) the Blue River Dam validation thresholds; (4) the trans-standard conversion
loss policy and signed \texttt{LossReport} schema; and (5) the academic foundation
citation. The 13 core standards are XES (IEEE 1849), OCEL 2.0, POWL, OCPQ, BPMN 2.0,
Petri Net, WF-Net, Declare, DFG, Process Tree, PROV-O, OTel Weaver, and DCTERMS. The
3 metadata standards are SHACL, SKOS, and ODRL.

\paragraph{The canonical process model $W_{\text{std}}$.}
Defined in \texttt{reverse-lock-in.md} as a WF-net:
\begin{equation}
  W_{\text{std}} = (P_{\text{std}}, T_{\text{std}}, F_{\text{std}}, i_{\text{std}}, o_{\text{std}})
\end{equation}
where $P_{\text{std}}$ is the set of authorized operational states (places), $T_{\text{std}}$
is the set of lawful transitions, $F_{\text{std}}$ is the set of legal flow relations, and
$i_{\text{std}}$, $o_{\text{std}}$ are the unique source and sink places. This model is
sound by WF-net soundness criteria, cryptographically signed and referenced by BLAKE3 hash,
and publicly archived. It is the formal referent of the reverse lock-in doctrine.

\paragraph{BLAKE3 receipt chain.}
The cryptographic primitive is BLAKE3 hash chaining. For XES traces:
$\mathcal{H}(\sigma) = \operatorname{BLAKE3}(\mathcal{H}(e_1) \parallel \cdots \parallel
\mathcal{H}(e_n))$. For the Slide-to-Receipt bridge:
$\text{SlideBinding} = \operatorname{BLAKE3}(\text{SlideUUID} \parallel \text{ReceiptHash}
\parallel \text{ClaimValue})$. For competitor proof requirements:
$\Pi = \operatorname{BLAKE3}(\text{predecessor\_receipt} \parallel \text{transition\_data}
\parallel \text{timestamp})$.

\paragraph{LossReport.}
Every trans-standard conversion (e.g., XES to OCEL 2.0, OCEL 2.0 to XES) emits a named,
signed artifact: the \texttt{LossReport}. It records the source format, target format,
structural changes (synthetic objects materialized, pruned attributes, discarded O2O
links), and an ED25519 witness signature. The existence of a signed \texttt{LossReport}
is a type-law requirement for any conversion; silent behavioral loss is architecturally
refused.

\subsection{Admissible Transitions}
\label{sec:standards-arch-transitions}

The Blue River Dam (\texttt{public\_standards\_to\_blue\_river\_dam.md}) defines three
mandatory validation gates that any process state transition must pass:

\begin{enumerate}
  \item \textbf{Fitness gate:} Aggregate replay fitness $f \ge 0.95$. Under no
        circumstances shall a transaction be admitted if fitness falls below $0.85$,
        even with a board override.
  \item \textbf{Precision gate:} Process model precision $p \ge 0.90$, preventing
        under-specified models from hiding deviations.
  \item \textbf{Boundedness gate:} The execution core automatically freezes and aborts
        any transaction that violates 1-boundedness of control places.
\end{enumerate}

The POWL Projection Law (v30.1.1 Spec) adds a fourth condition: a POWL model $M$ is
process-tree projectable if and only if $\mathcal{L}(M) = \mathcal{L}(T)$ for some
block-structured process tree $T$. If $M$ contains irreducible partial orders, it exceeds
the process tree language. The \texttt{TreeProjectable} trait is sealed and implemented
only for \texttt{ProcessTreeProjectable}; non-projectable POWL models cannot be silently
projected.

\subsection{Boundaries}
\label{sec:standards-arch-boundaries}

The standards project defines three explicit graduation boundary items that are
documented as open questions rather than completed capabilities:

\begin{itemize}
  \item \textbf{OCPQ query evaluation engine:} Structural shapes exist in
        \texttt{src/ocpq.rs} in the wasm4pm codebase, but the query evaluation engine
        itself has not graduated to wasm4pm. The OCPQ placement file is VERIFIED in the
        audit matrix, but the execution capability is deferred.
  \item \textbf{SHACL shapes for OCEL 2.0 RDF:} The SHACL/OCEL surface is described in
        the placement file, but no \texttt{.shacl} or \texttt{.ttl} files exist in the
        \texttt{standards/} directory.
  \item \textbf{OTel Weaver process-specific conventions:} The Weaver framework is mapped
        but the actual \texttt{.tera} or \texttt{.yaml} convention files are not present
        in \texttt{standards/}.
\end{itemize}

These boundaries are documented per the foundry's immutability doctrine: gaps are named,
not hidden.

\subsection{Failure Modes Refused}
\label{sec:standards-arch-failures}

The architecture explicitly refuses four failure modes:

\begin{enumerate}
  \item \textbf{Process log laundering:} Refused by XES trace hash chaining and the
        chronological invariance check $\forall i,\ t(e_{i+1}) \ge t(e_i)$.
  \item \textbf{Phantom object injection:} Refused by OCEL 2.0 structural validation
        via \texttt{OcelLog::validate}, which emits
        \texttt{OcelRefusal::DanglingEventObjectLink} for any event-to-object link whose
        object identifier does not appear in $\mathcal{O}$.
  \item \textbf{Silent behavioral loss in trans-standard conversion:} Refused by the
        mandatory \texttt{LossReport} requirement. Every conversion that discards
        structural information must emit a signed artifact naming the loss.
  \item \textbf{Unsound process model deployment:} Refused by WF-net soundness
        verification via the Karp-Miller Coverability Graph and Short-Circuited Petri Net
        analysis, as specified in \texttt{wf-net\_verification\_specification.md} v30.1.2.
\end{enumerate}
